% !TEX root = ../Presentation.tex
% ---------------------------------------------------------------------------
% Model comparison and recovery
% Slides taken from ../Presentation_ABC/sections; trimmed ones are marked.
% ---------------------------------------------------------------------------
\section{Model Comparison and Recovery}

\subsection{Predictive Performance}

% trimmed; numbers from Model_B/Fit_B_Report.tex and Model_C/Fit_C_Report.tex
\begin{frame}{Leave-one-out cross-validation}
  How well each model predicts each trial when that trial is left out of the fit
  (higher is better):
  \vspace{0.8em}

  \centering
  \begin{tabular}{lc}
    \toprule
    Model & difference from B (SE) \\
    \midrule
    A: sessions pooled          & $-164.1$ (19.2) \\
    B: sessions hierarchical    & \textbf{best} \\
    C: distractor as covariate  & $-37.1$ (11.9) \\
    \bottomrule
  \end{tabular}

  \vspace{0.8em}
  \raggedright
  \begin{itemize}
    \item B beats A clearly, despite many more parameters
    \item B beats C, although C needs fewer parameters
  \end{itemize}
\end{frame}

\subsection{Model Recovery: B vs A}
% numbers from Model_B/Model_Recov_BA_Report.tex

% trimmed
\begin{frame}{Model recovery: LOO picks the generating model}
  Is B's advantage real, or just flexibility? Simulate datasets from each model on the real
  design, fit both, and see which one LOO picks:
  \vspace{0.6em}

  \centering
  \begin{tabular}{lcc}
    \toprule
     & selects A & selects B \\
    \midrule
    data from A & 9 & 1 \\
    data from B & 0 & 10 \\
    \bottomrule
  \end{tabular}

  \vspace{0.6em}
  \raggedright
  \begin{itemize}
    \item On data from A, B ties or loses: flexibility alone does not make it win
    \item The real data look like the data from B: sessions really differ
  \end{itemize}
\end{frame}

\plotframe{Model recovery: $\Delta$ in every replicate}
  {\plotsB/Model_Recov_Plots_BA/elpd_diff_by_replicate.png}

\subsection{Model C vs Model B}
% numbers from Model_C/Fit_C_Report.tex, "Comparison with Model B"

\plotframe{Model C vs Model B: session PSE shifts}
  {\plotsC/Fit_Plots_C/pse_shift_C_vs_B.png}

% trimmed
\begin{frame}{Model comparison: conclusions}
  \begin{itemize}
    \item \textbf{Sessions differ:} B predicts far better than A, and model recovery shows this
          is not B's flexibility alone
    \item \textbf{The models agree on the PSE shift:} about 4~$\mu m$ (left) and 6~$\mu m$
          (right) on average
    \item \textbf{B predicts better than C mainly because of lapses:} C's one lapse rate for
          both conditions misses the extra ``high'' responses to low stimulus amplitudes on bimanual
          trials
    \item Less clearly: a distractor-presence effect (right hand) and a bimanual slope
    \item[$\Rightarrow$] Model B is the model to use for inference for now
  \end{itemize}
\end{frame}
